% ===== tr-08 · b15 · tr08e：第 5 章 p38--40（原文页 38--40）=====
% [原文 p.38]
\begin{proof}[证明（概要）]
\begin{itemize}
\item 我们首先用某个 NA Fubini--Study 度量在 $C^0$ 意义下逼近 NA 度量 $\|\cdot\|_{CY}$；该度量自然地作为 $X_t$ 上通常 Fubini--Study 度量在混合拓扑下的极限出现（参看第 5.5.3 节）。Fubini--Study 度量是正的，并且按构造其局部势与 $\phi_0\circ\mathrm{Log}_{\mathcal X}$ 在 $C^0$ 意义下相差任意小。
\item 我们无法直接控制 Fubini--Study 度量的体积形式；相应射影嵌入的次数巨大。相反，由第 5.3.1 节的体积渐近与实 MA 方程 (5.17)，局部势 $\phi_0\circ\mathrm{Log}_{\mathcal X}$ 的体积形式与 $\frac{(L^n)}{|\log|t||^n}d\mu_t$ 之差可忽略。
\item 想法是再做一次正则化。我们在 $X_t$ 的一般区域修改 Fubini--Study 度量，使其在一般区域与 $\phi_0\circ\mathrm{Log}_{\mathcal X}$ 相差 $C^2$-小误差。这一步我们也借助了实 MA 方程的正则性理论。最终得到 $\omega_{\psi,t}$，它按构造是 Kähler 的。
\item 在非一般区域，我们不作正则化。由于当 $|t|\ll 1$ 时一般区域已占 $\omega_{\psi,t}^n$ 测度的 $99.9\%$，非一般区域的总测度可忽略。我们用它来论证总变差界。\proofend
\end{itemize}
\end{proof}

% [Tutor Note: 本证明的主角是「先逼近、再正则化」两步：逼近给 $C^0$ 控制（但体积形式不可控），正则化才给出体积形式控制；非一般区域根本不作正则化，只靠其测度可忽略（$99.9\%$ 是作者给出的启发式数值，原文并未声称这是定理）。注意体积形式的比较对象是 $\frac{(L^n)}{|\log|t||^n}d\mu_t$，其归一化因子 $(L^n)$ 来自相交数。】

\subsubsection{势的 $C^0$-收敛}
还需证明

\begin{proposition}{5.6.5}
在把区域略微收缩之后，$U_{J,t}$ 上的 Calabi--Yau 度量容许局部势函数 $\phi_{CY,J,t}$，使得 $\omega_{CY,t}=dd^c\phi_{CY,J,t}$，且当 $t\to 0$ 时 $\|\phi_{CY,J,t}-\phi_0\circ\mathrm{Log}_{\mathcal X}\|_{C^0}\to 0$。
\end{proposition}

命题 5.6.4 说明，在 $t\to 0$ 极限下 $\omega_{\psi,t}$ 的局部势与 $\phi_0\circ\mathrm{Log}_{\mathcal X}$ 在 $C^0$ 意义下相差可忽略。理想情况下，人们希望用某个版本的 $L^1$-体积稳定性来推出 $\omega_{\psi,t}$ 与 $\omega_{CY,t}$ 之间相对势的 $C^0$-小性。不幸的是，由于非一般区域中正则化的困难，我们对 $\omega_{\psi,t}$ 在该区域的体积密度几乎没有控制，无法对 $\omega_{\psi,t}$ 得出像式 (5.11) 那样的 Skoda 型估计。这一技术问题造成 $\omega_{\psi,t}$ 与 $\omega_{CY,t}$ 之间的不对称，只有 $L^1$-体积稳定性估计（参看第 5.4.3 节）的「一半」适用；这正是我们设计定理 5.4.7 的原因。

尘埃落定之后，定理 5.4.7 蕴含 $\phi_{CY,J,t}-\phi_0\circ\mathrm{Log}_{\mathcal X}$ 在一个几乎占 Calabi--Yau 测度 $100\%$ 的子集上集中在其最小值（归一化为零）附近（参看 [53, Prop 4.4, Cor. 4.6]）。

% [Tutor Note: 命题 6.5 是「假设」而非「结论」：它要求的正是 $C^0$ 收敛，用来喂给命题 6.3 以得到 $C^\infty$-渐近。原文说「只能用到 $L^1$-体积稳定性的一半」，原因是 $\omega_{\psi,t}$ 在非一般区域的体积密度不受控，故 Skoda 型估计只能单侧使用------这也是定理 4.7 的设计动机。]

% [原文 p.39]
更精确地说，对任意给定的小数 $\kappa,\lambda\ll 1$，当 $t$ 充分小时，有
\begin{equation}
d\mu_t(\phi_{CY,J,t}-\phi_0\circ\mathrm{Log}_{\mathcal X}\ge\kappa/4)<\lambda. \tag{5.21}
\end{equation}
在 $U_{J,t}$ 略微收缩后的版本上，这可以改进为 $C^0$-控制
\[
0\le\phi_{CY,J,t}-\phi_0\circ\mathrm{Log}_{\mathcal X}<\kappa,
\]
这要用中值不等式作一次颇为棘手的应用（参看 [53, Thm. 4.7]）。由于 $\kappa$ 任意，这就实现了命题 5.6.5，从而验证了命题 5.6.5 的假设，于是得到 SYZ 猜想的弱度量版本。

% [Tutor Note: (21) 的门槛是 $\kappa/4$（不是 $\kappa/2$ 或 $\kappa$），这个具体常数来自把测度估计与中值不等式拼接时的技术性因子，不要擅自「化简」。本段末尾「实现命题 6.5，从而验证了命题 6.5 的假设」原文两处都写作 Prop. 6.5（p39，行 "Since $\kappa$ is arbitrary..." 与下一行起首），按逻辑第二处应为命题 6.3（其假设正是一致 $C^0$-收敛）；译文照原文排并在此注明。]

\begin{remark}{5.16}
上述 $C^0$ 收敛的陈述只适用于一般区域。
\end{remark}

我们不知道下述问题的答案：

\begin{question}{5.7}
当 $t\to 0$ 时，$X_t$ 上 CY 度量的势是否在混合拓扑下整体地收敛到 $X_K^{an}$ 上的 NA CY 度量 $\|\cdot\|_{CY}$？
\end{question}

% [Tutor Note: 注 16 与问题 7 是一对：注 16 承认 $C^0$ 收敛只在一般区域成立（非一般区域靠测度可忽略绕过去），问题 7 则问整体收敛是否成立------这正是下一节「策略 II」要正面处理的问题。问题编号原文为 Question 7，本章内套用 $\{5.7\}$。]

\subsection{策略 II：先验极限}
第二个策略不借助 NA 几何，且与第 5.6.2 节无关。
\begin{itemize}
\item 先验地论证：$U_{J,t}\subset X_t$ 上 Calabi--Yau 度量的局部势函数在 $C^0$-范数下子序列收敛到 $Sk(X)$ 的开 $n$ 维面上的某个凸函数；
\item 再论证该凸函数满足实 MA 方程。
\end{itemize}
这一策略具有一般性，只除开一个微妙的问题------它只在 Fermat 族这一非常特殊的情形下被我们解决（参看定理 5.6.2 [52]）。

\subsubsection{构造凸函数}
回忆对数映射 $\mathrm{Log}_t:(\C^*)^n\to\R^n$：
\[
\mathrm{Log}_t(z_1,\dots z_n)=\frac{1}{\log|t|}(\log|z_1|,\dots\log|z_n|).
\]
考虑开凸子集 $U\subset\R^n$，并设 $\phi$ 是 $\mathrm{Log}_t^{-1}(U)\subset(\C^*)^n$ 上的 psh 函数。

\begin{lemma}{5.6.6}[52, Lemma 4.3]
逐纤维的 $T^n$ 平均函数
\[
\bar\phi(x_1,\dots x_n)=\frac{1}{(2\pi)^n}\int_{T^n}\phi(e^{x_1\log|t|+i\theta_1},\dots e^{x_n\log|t|+i\theta_n})d\theta_1\dots d\theta_n
\]
作为 $x_1,\dots x_n$ 的函数是凸的。
\end{lemma}

% [Tutor Note: 这里出现两个不同的对数映射，务必区分：$\mathrm{Log}_{\mathcal X}:X_t\to\Delta_{\mathcal X}$ 依赖 snc 模型（第 6.1 节），而本页的 $\mathrm{Log}_t:(\C^*)^n\to\R^n$ 是局部坐标公式 $\frac{1}{\log|t|}(\log|z_1|,\dots)$。引理 6.6 的指数里是 $x_k\log|t|+i\theta_k$（$\log|t|$ 为实数因子、$i\theta$ 为虚部），不是 $\log(z_k)$。$d\theta_1\dots d\theta_n$ 前面没有 $\frac{1}{(2\pi)^n}$ 之外的任何因子。]

% [原文 p.40]
\begin{proof}
由于函数 $\bar\phi$ 是 psh 函数的平均，它作为 $\mathrm{Log}_t^{-1}(U)$ 上的 $T^n$-不变函数是 psh 的。这类函数对应于底下的凸函数。\proofend
\end{proof}

一个重要的直观是：在 $(\C^*)^n$ 中充分坍缩的环面区域上，有界的 Kähler 势强烈倾向于被凸函数逼近。

\begin{proposition}{5.6.7}
假设 $\|\phi\|_{C^0}$ 具有与 $t$ 无关的一致界。则把 $U$ 收缩一个与 $t$ 无关的小量之后，有
\begin{itemize}
\item 凸函数 $\bar\phi$ 具有 Lipschitz 界 $|\bar\phi(x)-\bar\phi(x')|\le C|x-x'|$；
\item 有上界 $\phi-\bar\phi\le\frac{C}{|\log|t||^{1/2}}$；
\item 在每个对数二进尺度 $U_a=\{a_i\le\log|z_i|\le 2a_i,\ \forall i\}\subset U$ 上，$L^1$-积分
\[
\int_{U_a}|\phi-\bar\phi|\prod\sqrt{-1}d\log z_i\wedge d\,\overline{\log z_i}\le\frac{C}{|\log|t||^{1/2}};
\]
\item 存在改进的 Skoda 不等式，其常数 $\alpha,C$ 一致且与 $t$ 无关：
\[
\int_U e^{-\alpha|\log|t||^{1/2}(\phi-\bar\phi)}d\mu_t\le C.
\]
\end{itemize}
\end{proposition}

\begin{proof}[证明（概要）]
\begin{itemize}
\item 有界凸函数在略微收缩的凸区域上自动具有 Lipschitz 界。
\item 第二条来自对次调和函数的中值不等式作一次颇为棘手的应用，参看 [52, section 4.3]。
\item 第三条是因为 $\phi-\bar\phi$ 的均值为零，故上界蕴含 $L^1$-界，参看 [52, section 4.3]。
\item 先对函数 $\phi$ 在每个对数二进尺度上应用基本 Skoda 估计（定理 5.4.2），在这些尺度上由 Lipschitz 界 $\bar\phi$ 几乎是常数；然后对所有对数二进尺度求和（参看 [52, section 4.6]）。\proofend
\end{itemize}
\end{proof}

% [Tutor Note: 命题 6.7 第三条的被积式是 $|\phi-\bar\phi|$ 乘上 $(\C^*)^n$ 的体积形式 $\prod\sqrt{-1}d\log z_i\wedge d\overline{\log z_i}$（第二个因子原文的横线覆盖整个 $\log z_i$，与 (19) 的写法一致）；$U_a=\{a_i\le\log|z_i|\le 2a_i\}$ 是「对数二进」盒子，注意里面的绝对值是加在 $z_i$ 上（$\log|z_i|$），不是 $|\log z_i|$。第四条指数中 $|\log|t||^{1/2}$ 是外层绝对值再开方，与第二条的分母同为 $1/2$ 次幂------正是这个 $1/2$ 使衰减比通常 Skoda 估计更强。]

\begin{remark}{5.17}
改进的 Skoda 估计是 [52] 中的主要发现之一。直观地说，这意味着 $\phi-\bar\phi$ 只能在测度指数小的集合上显著低于 $-\frac{\mathrm{Const}}{|\log|t||^{1/2}}$。与上界 $\phi-\bar\phi\le\frac{C}{|\log|t||^{1/2}}$ 合起来，这意味着当 $t$ 充分小时，任意有界 psh 函数 $\phi$ 除在一个指数小测度的集合上之外，都非常接近凸函数 $\bar\phi$。
\end{remark}

% [Tutor Note: 注 17 的论证是「单侧 Skoda + 上界」夹逼：下界只在指数小测度集上被破坏，而上界处处成立，故 $\phi$ 与凸函数 $\bar\phi$ 的差被控制在 $|\log|t||^{-1/2}$ 量级。$\mathrm{Const}$ 在原文中排成正体（不是一个具体的 $C$），与同段的上界常数 $C$ 不同，译文保留正体。]
